\section{设计过程与实施}
\subsection{资料调研与方案演化}
记录素材来源、授权范围及设计约束，用草图、迭代记录和最终作品说明设计决策。官方艺术类文件明确允许说明书版式自定，因此本节是可修改的设计报告结构，不是学校规定的章节清单。
\subsection{作品图版示例}
以下为矢量构图占位示例，不是毕业设计作品；请替换为具有使用权的作品照片或效果图。
\begin{figure}[htbp]\centering
\begin{tikzpicture}
\fill[black!8] (0,0) rectangle (9,4);
\draw[very thick] (0.6,0.6) rectangle (4,3.4);
\draw[very thick] (5,0.6)--(8.4,3.4) (5,3.4)--(8.4,0.6);
\node at (4.5,-0.4) {图版占位：替换为实际作品};
\end{tikzpicture}
\caption{构图占位示例（非作品成果）}
\end{figure}
\section{反思与成果归档}
说明设计目标的达成情况、制作成本、材料限制、改进方向和作品版权。按官方附件要求另备 A4 彩色作品图版；研究涉及非遗图案、人物肖像或第三方素材时，应注明出处并核对授权。
